% ch03_u04 —— 原书 p.91–p.100（3.6.1 末、3.6.2–3.6.19、3.7 开头至 3.7.5）
\SourcePage{91}

其中 $z_{1}, z_{2} \in \mathfrak{Z}$，而 $\tilde{z}_{i}$ 是 $z_{i}$ 在 $A$ 中的原像。Hayashi 在 [Ha88] 中指出：即使去掉假设 $\mathfrak{Z} = A_{0}$，（79）仍是 $\mathfrak{Z}$ 上的良定义泊松括号。特别地，（79）的右端属于 $\mathfrak{Z}$。

\emph{注记}
\begin{enumerate}[label=(\roman*),leftmargin=2.2em]
  \item 上述情形有典范的李代数态射 $\varphi: \mathfrak{Z} \to \operatorname{Der} A_{0}/\operatorname{Int} A_{0}$，其中 $\operatorname{Int} A_{0}$ 是内导子构成的空间。$\varphi$ 由 $\varphi(z) = D_{z}$ 定义，其中 $D_{z}(a) := [\tilde{z}, \tilde{a}]/h \bmod h$，而 $\tilde{z}, \tilde{a} \in A$ 分别是 $z \in \mathfrak{Z}$ 与 $a \in A_{0}$ 的原像。若 $z' \in \mathfrak{Z}$，则 $D_{z}(z') = \{z, z'\}$。$\operatorname{Der} A_{0}/\operatorname{Int} A_{0}$ 是 $\mathfrak{Z}$-模，且 $\varphi(z_{1}z_{2}) = z_{1}\varphi(z_{2}) + z_{2}\varphi(z_{1})$。于是 $\varphi$ 诱导 $\mathfrak{Z}$-模态射 $\Phi: \Omega^{1}_{\mathfrak{Z}} \to \operatorname{Der} A_{0}/\operatorname{Int} A_{0}$。事实上 $\Phi$ 是 $\mathfrak{Z}$ 上李代数胚的态射（李代数胚的定义见 3.5.10）。$\operatorname{Der} A_{0}/\operatorname{Int} A_{0}$ 上的李代数胚结构按显然方式定义；$\Omega^{1}_{\mathfrak{Z}}$ 上的结构则是泊松括号诱导的标准代数胚结构（参看 [We88]）。即对 $z, z' \in \mathfrak{Z}$ 令 $[dz, dz'] := d\{z, z'\}$；而态射 $\Omega^{1}_{\mathfrak{Z}} \to \operatorname{Der} \mathfrak{Z}$ 把 $dz$ 映到 $\operatorname{grad} z$，$(\operatorname{grad} z)(z') := \{z, z'\}$。
  \item 若把 $\mathbb{C}[h]$ 换成 $\mathbb{C}[h]/(h^{3})$，上述构造仍有意义。
\end{enumerate}

\bdp{3.6.2} 现在设 $\mathfrak{g}$ 是半单李代数，$K := \mathbb{C}((t))$。记 $A$ 为 2.5.1 中李代数 $\widetilde{\mathfrak{g}\otimes K}$ 的完备化泛包络代数，即 $A := \varprojlim_{n}(U\widetilde{\mathfrak{g}\otimes K})/J_{n}$，其中 $J_{n} \subset U\widetilde{\mathfrak{g}\otimes K}$ 是由 $\mathfrak{g}\otimes t^{n}\mathbb{C}[[t]] \subset \mathfrak{g}\otimes K \subset \widetilde{\mathfrak{g}\otimes K}$ 生成的左理想，$n \ge 0$。考虑 $A$ 上由 $ha = \mathbf{1}\cdot a - a$（$a \in A$）定义的 $\mathbb{C}[h]$-代数结构，其中 $\mathbf{1} \in \mathbb{C} \subset \widetilde{\mathfrak{g}\otimes K} \subset A$。$A$ 在 $\mathbb{C}[h]$ 上平坦，而 $A/hA$ 是 2.5.2 与 2.9.4 中的完备化扭泛包络代数 $\overline{U}' = \overline{U}'(\mathfrak{g}\otimes K)$。于是（79）定义了 $\overline{U}'$ 的中心 $\mathfrak{Z}$ 上的泊松括号。[Ha88] 引入了这个括号，故称之为 Hayashi 括号。

\bdp{3.6.3} 对开的李子代数 $\mathfrak{a} \subset \mathfrak{g}\otimes O$，记 $\mathcal{J}_{\mathfrak{a}}$（相应地 $\widetilde{\mathcal{J}}_{\mathfrak{a}}$）为 $\overline{U}'$（相应地 $A = \overline{U}\widetilde{\mathfrak{g}\otimes K}$）中由 $\mathfrak{a} \subset \mathfrak{g}\otimes O \subset \mathfrak{g}\otimes K$ 生成的左理想的闭包。显然 $\mathcal{J}_{\mathfrak{a}}$ 是 $\widetilde{\mathcal{J}}_{\mathfrak{a}}$ 在 $\overline{U}'$ 中的像。令 $I_{\mathfrak{a}} := \mathcal{J}_{\mathfrak{a}} \cap \mathfrak{Z}$。我们给 $\mathfrak{Z}$ 装上由 $\overline{U}'$ 诱导的拓扑。

\SourcePage{92}

理想 $\mathcal{J}_{\mathfrak{a}}$（相应地 $I_{\mathfrak{a}}$）构成 $\overline{U}'$（相应地 $\mathfrak{Z}$）中零的邻域基。

\bdp{3.6.4} \emph{引理.}
\begin{enumerate}[label=(\roman*),leftmargin=2.2em]
  \item $\{I_{\mathfrak{a}}, I_{\mathfrak{a}}\} \subset I_{\mathfrak{a}}$。
  \item $\mathfrak{Z}$ 上的 Hayashi 括号连续。
\end{enumerate}

\emph{证明.} 用到如下事实：$A/\widetilde{\mathcal{J}}_{\mathfrak{a}}$ 带 3.6.2 的 $\mathbb{C}[h]$-模结构时平坦。\hfill$\Box$

\bdp{3.6.5} 令 $I := I_{\mathfrak{g}\otimes O}$。典范态射 $\mathfrak{Z} \to \mathfrak{z}_{\mathfrak{g}}(O)$ 是满射（见 2.9.3--2.9.5），其核等于 $I$。于是 $\mathfrak{z}_{\mathfrak{g}}(O) = \mathfrak{Z}/I$。

记 $I^{2}$ 为 $\mathfrak{Z}$ 中由形如 $ab$（$a, b \in I$）的元素生成的闭理想。则 $I/I^{2}$ 是 $\mathfrak{z}_{\mathfrak{g}}(O)$ 上的李代数胚：交换子 $I/I^{2} \times I/I^{2} \to I/I^{2}$ 与映射 $I/I^{2} \to \operatorname{Der}\mathfrak{z}_{\mathfrak{g}}(O)$ 都由 Hayashi 括号诱导。李代数 $\operatorname{Der} O$ 作用在 $I/I^{2}$ 与 $\mathfrak{z}_{\mathfrak{g}}(O)$ 上。这些作用连续：$I/I^{2}$ 带由 $\mathfrak{Z}$ 诱导的拓扑，$\mathfrak{z}_{\mathfrak{g}}(O)$ 取离散拓扑。

\bdp{3.6.6} 我们来表述定理 3.2.2 的一个更精确的版本。已有代数 $\mathfrak{z}_{\mathfrak{g}}(O)$ 及其上的李代数胚 $I/I^{2}$。另一方面，记 ${}^{L}\mathfrak{g}$ 为朗兰兹对偶，考虑代数 $A_{{}^{L}\mathfrak{g}}(O)$（见 3.2.1）及其上的李代数胚 $\mathfrak{a}_{{}^{L}\mathfrak{g}}$（见 3.5.11）。$I/I^{2}$ 与 $\mathfrak{a}_{{}^{L}\mathfrak{g}}$ 都带拓扑（见 3.6.5 与 3.5.15）。李代数 $\operatorname{Der} O$ 作用在所有这些对象上。$\mathfrak{z}_{\mathfrak{g}}(O)$ 与 $A_{{}^{L}\mathfrak{g}}(O)$ 都带滤链（见 1.2.5 与 3.2.1）。此外有态射 $\sigma_{A}^{-1}\sigma_{\mathfrak{z}}: \operatorname{gr}\mathfrak{z}_{\mathfrak{g}}(O) \to \operatorname{gr} A_{{}^{L}\mathfrak{g}}(O)$，其中 $\sigma_{\mathfrak{z}}: \operatorname{gr}\mathfrak{z}_{\mathfrak{g}}(O) \to \mathfrak{z}^{cl}_{\mathfrak{g}}(O)$ 是象征映射，而 $\sigma_{A}$ 是把 $\mathfrak{g}$ 换成 ${}^{L}\mathfrak{g}$ 后的同构（48）。

\bdp{3.6.7} \emph{定理.} 存在滤 $\operatorname{Der} O$-代数的同构
\begin{equation}\tag{80}
\varphi_{O}: A_{{}^{L}\mathfrak{g}}(O) \xrightarrow{\sim} \mathfrak{z}_{\mathfrak{g}}(O)
\end{equation}

\SourcePage{93}
使得 $\operatorname{gr}\varphi_{O}^{-1} = \sigma_{A}^{-1}\sigma_{\mathfrak{z}}$；并且 $\varphi_{O}$ 可延拓为李代数胚的拓扑 $\operatorname{Der} O$-等变同构
\begin{equation}\tag{81}
\mathfrak{a}_{{}^{L}\mathfrak{g}} \xrightarrow{\sim} I/I^{2}.
\end{equation}

本定理可从 [FF92] 中提取（见 3.7.12--3.7.17）。

\emph{注.} 由 3.5.14，同构（80）与（81）唯一。3.6.11 中我们将给出同构（81）的一条附加性质。但首先须为代数胚 $I/I^{2}$ 定义（78）的类比。

\bdp{3.6.8} 沿用 3.5.17 的记号。

\emph{引理.} 令 $\mathcal{J}_{-} := (\overline{U}')^{\le 0} \cap \mathcal{J}_{\mathfrak{a}}$，其中 $\mathfrak{a} = t\mathfrak{g}[[t]]$，而 $\mathcal{J}_{\mathfrak{a}}$ 见 3.6.3。则 $\mathcal{J}_{-}$ 是 $(\overline{U}')^{\le 0}$ 的双边理想，并且
\begin{equation}\tag{82}
(\overline{U}')^{\le 0} = U\mathfrak{g} \oplus \mathcal{J}_{-}.
\end{equation}

\emph{证明.} （82）显然。由于 $\mathcal{J}_{-}$ 是左理想且 $[\mathfrak{g}, \mathcal{J}_{-}] \subset \mathcal{J}_{-}$，（82）表明 $\mathcal{J}_{-}$ 是双边理想。\hfill$\Box$

定义 $\pi: (\overline{U}')^{\le 0} \to U\mathfrak{g}$ 为满足 $\pi(\mathcal{J}_{-}) = 0$ 且 $\pi(a) = a$（$a \in U\mathfrak{g}$）的态射。

这里给出 $\pi$ 的等价定义。令 $\operatorname{Vac}'_{\mathfrak{a}} := \overline{U}'/\mathcal{J}_{\mathfrak{a}}$，$\mathfrak{a} = t\mathfrak{g}[[t]]$。则 $\operatorname{Vac}'_{\mathfrak{a}}$ 是左 $\overline{U}'$-模，也是右 $U\mathfrak{g}$-模。$L_{0}$ 在 $\operatorname{Vac}'_{\mathfrak{a}}$ 上的特征值非负，且 $\operatorname{Ker}(L_{0}: \operatorname{Vac}'_{\mathfrak{a}} \to \operatorname{Vac}'_{\mathfrak{a}}) = U\mathfrak{g}$。于是 $U\mathfrak{g} \subset \operatorname{Vac}'_{\mathfrak{a}}$ 在 $(\overline{U}')^{\le 0}$ 的左作用下不变。$(\overline{U}')^{\le 0}$ 的左作用与 $U\mathfrak{g}$ 的右作用交换，故它定义了态射 $(\overline{U}')^{\le 0} \to U\mathfrak{g}$。这就是 $\pi$。

\bdp{3.6.9} 记 $C$ 为 $U\mathfrak{g}$ 的中心。则
\[
  \pi(\mathfrak{Z}^{\le 0}) \subset C, \qquad \pi(\mathfrak{Z}\cdot\mathfrak{Z}^{<0} \cap \mathfrak{Z}^{\le 0}) = 0.
\]
设 $m \subset C$ 是 $U\mathfrak{g}$ 的单位表示对应的极大理想。回顾 $I := \operatorname{Ker}(\mathfrak{Z} \to \mathfrak{z}_{\mathfrak{g}}(O))$。则 $\pi(I^{\le 0}) \subset m$。由于 $(I^{2})^{\le 0} \subset I^{\le 0}\cdot I^{\le 0} + (\mathfrak{Z}\cdot\mathfrak{Z}^{<0} \cap \mathfrak{Z}^{\le 0})$，有 $\pi((I^{2})^{\le 0}) \subset m^{2}$。于是 $\pi$ 诱导 $\mathbb{C}$-线性映射 $d: (I/I^{2})^{\le 0} \to m/m^{2}$，使得 $\mathfrak{z}_{\mathfrak{g}}(O)\cdot(I/I^{2})^{<0} \cap (I/I^{2})^{\le 0} \subset \operatorname{Ker} d$。

\SourcePage{94}

\emph{练习.} 对 $z_{1}, z_{2} \in \mathfrak{Z}^{\le 0}$ 有 $\pi(\{z_{1}, z_{2}\}) = 0$（于是 $d$ 是李代数态射）。

\bdp{3.6.10} 把 $C$ 与 $\mathfrak{h}^{*}$ 上 $W$-不变多项式的代数等同，其中 $W$ 是 Weyl 群。则 $m$ 由 $\mathfrak{h}^{*}$ 上在 $\rho$ 处为零的 $W$-不变多项式构成，其中 $\rho :=$ 基本权之和。由于 $\rho \in \mathfrak{h}^{*}$ 是正则的，可把 $m/m^{2}$ 与 $\mathfrak{h}$ 等同：把 $m$ 中的 $W$-不变多项式对应到它在 $\rho$ 处的微分。于是我们构造了映射
\begin{equation}\tag{83}
d: (I/I^{2})^{\le 0}/\bigl(\mathfrak{z}_{\mathfrak{g}}(O)\cdot(I/I^{2})^{<0} \cap (I/I^{2})^{\le 0}\bigr) \to \mathfrak{h}
\end{equation}

\bdp{3.6.11} \emph{定理.} 下图反交换：
\begin{equation}\tag{84}
\begin{array}{ccc}
(\mathfrak{a}_{{}^{L}\mathfrak{g}})^{\le 0}/\bigl(A_{{}^{L}\mathfrak{g}}(O)\cdot\mathfrak{a}_{{}^{L}\mathfrak{g}}^{<0} \cap (\mathfrak{a}_{{}^{L}\mathfrak{g}})^{\le 0}\bigr) & \xrightarrow{\ \sim\ } & \mathfrak{h}^{*}\\[2mm]
\big\downarrow\ \cong & & \big\uparrow\ \cong\\[2mm]
d: (I/I^{2})^{\le 0}/\bigl(\mathfrak{z}_{\mathfrak{g}}(O)\cdot(I/I^{2})^{<0} \cap (I/I^{2})^{\le 0}\bigr) & \longrightarrow & \mathfrak{h}
\end{array}
\end{equation}
这里上方箭头是把 $\mathfrak{g}$ 换成 ${}^{L}\mathfrak{g}$ 后的同构（78），左方箭头由（81）诱导，右方箭头来自数积（18）。

本定理可从 [FF92] 中提取（见 3.8.15--3.8.22）。

\bdp{3.6.12} 「临界」数积（18）在 3.6.11 中出现，原因并不严重。读者或许更喜欢下述观点。记 $B$ 为 $\mathfrak{g}$ 上不变双线性型的集合。对每个 $b \in B$，有完备化扭泛包络代数 $\overline{U}'_{b} = \overline{U}'_{b}(\mathfrak{g}\otimes K)$，它对应余圈 $(u, v) \mapsto \operatorname{Res} b(du, v)$，$u, v \in \mathfrak{g}\otimes K$（故 $\overline{U}' = \overline{U}'_{c}$，其中 $c$ 由（18）定义）。把 3.6.1 的一般构造施于依赖参数 $h$ 的代数族 $\overline{U}'_{c+hb}$，就得到与 $b \in B$ 对应的 $\mathfrak{Z}$ 上的泊松括号 $\{\}_{b}$（3.6.2 的括号对应 $b = c$）。$I/I^{2}$ 上的李代数胚结构依赖于 $b$。于是对每个非退化 $b \in B$，3.6.7 与 3.6.11 都成立（注意（84）中两个竖直箭头都依赖于 $b$）。

\SourcePage{95}

\bdp{3.6.13} 事实上，3.6.6--3.6.7 中提到的 $\operatorname{Der} O$ 在 $I/I^{2}$ 上的作用，来自典范态射 $\operatorname{Der} O \to I$，它本质上属于 Sugawara。在 3.6.14--3.6.15 简要回顾 Sugawara 公式之后，我们将在 3.6.16 解释这一点。这些公式还给出 $\mathfrak{z}_{\mathfrak{g}}(O)$ 的元素；当 $\mathfrak{g} = \mathfrak{sl}_{2}$ 时，它们生成 $\mathfrak{z}_{\mathfrak{g}}(O)$。我们在 3.6.18 中回顾这一点。3.6.18 与 3.6.19 在后文中都用不到 (?)。
% [Tutor Note: 原书 p.95 末括注「(?)」为作者自记之疑问，底本如此，保留原样。]

\bdp{3.6.14} 本小节回顾一般的 Sugawara 公式，3.6.15 回顾它们在临界水平的推论。

设 $A$ 为 $\widetilde{\mathfrak{g}\otimes K}$ 的完备化泛包络代数。作为向量空间，$\widetilde{\mathfrak{g}\otimes K}$ 是 $\mathfrak{g}\otimes K$ 与 $\mathbb{C} = \mathbb{C}\cdot 1$ 的直和。Sugawara 元素 $\widetilde{\mathcal{L}}_{n} \in A$ 由下式定义：
\begin{equation}\tag{85}
\widetilde{\mathcal{L}}_{n} := \frac{1}{2}\sum_{r+l=n} g^{\lambda\mu} :e_{\lambda}^{(r)}e_{\mu}^{(l)}:
\end{equation}
其中 $\{e_{\lambda}\}$ 是 $\mathfrak{g}$ 的一组基，$e_{\lambda}^{(r)} := e_{\lambda}t^{r} \in \mathfrak{g}((t)) = \mathfrak{g}\otimes K \subset \widetilde{\mathfrak{g}\otimes K}$，$(g^{\lambda\mu})$ 是 Gram 矩阵 $(e_{\lambda}, e_{\mu})$ 关于「临界」数积（18）的逆，并且
\begin{equation}\tag{86}
:e_{\lambda}^{(r)}e_{\mu}^{(l)}: \,=
\begin{cases}
e_{\lambda}^{(r)}e_{\mu}^{(l)} & \text{若 } r \le l,\\
e_{\mu}^{(l)}e_{\lambda}^{(r)} & \text{若 } r > l.
\end{cases}
\end{equation}
当然，（85）中对 $\lambda$ 与 $\mu$ 的求和是隐含的。显然无穷级数（85）收敛，且当 $n \to \infty$ 时 $\widetilde{\mathcal{L}}_{n} \to 0$。

\emph{注.} 若 $n \ne 0$，则（85）中的 $:e_{\lambda}^{(r)}e_{\mu}^{(l)}:$ 可换成 $e_{\lambda}^{(r)}e_{\mu}^{(l)}$。事实上，由于 $g^{\lambda\mu}$ 对称，除非 $r + l = 0$ 且 $r \ne 0$，总有 $g^{\lambda\mu}[e_{\lambda}^{(r)}, e_{\mu}^{(l)}] = 0$。

\SourcePage{96}

下述公式的证明可在文献中找到\footnote{读者应注意，Kac--Moody 代数的专家通常给 $\mathfrak{g}$ 配上把（18）除以负的对偶 Coxeter 数所得的数积。}，例如见 [KR] 的 Lecture 10 与 [Kac90] 的 §12.8：
\begin{equation}\tag{87}
\operatorname{ad}\widetilde{\mathcal{L}}_{n} = hL_{n},
\end{equation}
\begin{equation}\tag{88}
L_{m}(\widetilde{\mathcal{L}}_{n}) = (m-n)\widetilde{\mathcal{L}}_{m+n} + \delta_{m,-n}\cdot\frac{m^{3}-m}{12}\cdot(\dim\mathfrak{g})\cdot\mathbf{1}.
\end{equation}
（87）中 $\operatorname{ad}\widetilde{\mathcal{L}}_{n}$ 是算子 $A \to A$；$L_{n} := -t^{n+1}\frac{d}{dt} \in \operatorname{Der} K$ 也看作算子 $A \to A$（李代数 $\operatorname{Der} K$ 按显然方式作用在 $A$ 上）；$h$ 的含义同 3.6.2，即 $h: A \to A$ 是乘以 $\mathbf{1} - 1$。

利用（87）可把（88）改写成 Virasoro 形式：
\begin{equation}\tag{89}
[\widetilde{\mathcal{L}}_{m}, \widetilde{\mathcal{L}}_{n}] = h\Bigl((m-n)\widetilde{\mathcal{L}}_{m+n} + \delta_{m,-n}\cdot\frac{m^{3}-m}{12}\cdot(\dim\mathfrak{g})\cdot\mathbf{1}\Bigr).
\end{equation}

\bdp{3.6.15} $\widetilde{\mathcal{L}}_{n}$ 在 $A/hA = \overline{U}'$ 中的像记作 $\mathcal{L}_{n}$。由（87），$\mathcal{L}_{n}$ 属于中心 $\mathfrak{Z} \subset \overline{U}'$，并且
\begin{equation}\tag{90}
\{\mathcal{L}_{n}, z\} = L_{n}(z), \qquad z \in \mathfrak{Z}
\end{equation}
其中 $\{\ \}$ 表示 $\mathfrak{Z}$ 上的 Hayashi 括号。由（88）与（89）有
\begin{equation}\tag{91}
L_{m}(\mathcal{L}_{n}) = (m-n)\mathcal{L}_{m+n} + \delta_{m,-n}\cdot\frac{m^{3}-m}{12}\cdot\dim\mathfrak{g}
\end{equation}
\begin{equation}\tag{92}
\{\mathcal{L}_{m}, \mathcal{L}_{n}\} = (m-n)\mathcal{L}_{m+n} + \delta_{m,-n}\cdot\frac{m^{3}-m}{12}\cdot\dim\mathfrak{g}.
\end{equation}

\bdp{3.6.16} 若 $n \ge -1$，则 $\mathcal{L}_{n} \in I := \operatorname{Ker}(\mathfrak{Z} \to \mathfrak{z}_{\mathfrak{g}}(O))$。事实上，看（85）即知 $\mathcal{L}_{n}$ 消灭 $\operatorname{Vac}'$ 中的真空向量。若 $m, n \ge -1$，则「Virasoro 项」$\delta_{m,-n}(m^{3}-m)$ 为零，故有连续李代数态射 $\operatorname{Der} O \to I$，由 $L_{n} \mapsto \mathcal{L}_{n}$（$n \ge -1$）定义。它诱导连续代数态射
\begin{equation}\tag{93}
\operatorname{Der} O \to I/I^{2}.
\end{equation}

\SourcePage{97}

\emph{注.} 由（90），（93）诱导的 $\operatorname{Der} O$ 在 $I/I^{2}$ 上的作用与 3.6.6--3.6.7 中考虑的作用一致。

\bdp{3.6.17} \emph{引理.} （93）与同构 $I/I^{2} \xrightarrow{\sim} \mathfrak{a}_{{}^{L}\mathfrak{g}}$（它是（81）之逆）的复合，等于 3.5.11 中的态射 $\operatorname{Der} O \to \mathfrak{a}_{{}^{L}\mathfrak{g}}$。

\emph{证明.} 这两个态射 $\operatorname{Der} O \to \mathfrak{a}_{{}^{L}\mathfrak{g}}$ 诱导 $\operatorname{Der} O$ 在 $\mathfrak{a}_{{}^{L}\mathfrak{g}}$ 上的同一个作用。故由 3.5.12 它们相等。\hfill$\Box$

\bdp{3.6.18} 记 $\overline{\mathcal{L}}_{n}$ 为 $\mathcal{L}_{n}$ 在 $\mathfrak{Z}/I = \mathfrak{z}_{\mathfrak{g}}(O)$ 中的像。若 $n \ge -1$，则 $\overline{\mathcal{L}}_{n} = 0$。自然态射 $\mathbb{C}[\overline{\mathcal{L}}_{-2}, \overline{\mathcal{L}}_{-3}, \ldots] \to \mathfrak{z}_{\mathfrak{g}}(O)$ 是单射；当 $\mathfrak{g} = \mathfrak{sl}_{2}$ 时它是同构。为看出这一点，只需计算 $\overline{\mathcal{L}}_{n}$ 的主象征，并利用 2.4.1 对 $\mathfrak{z}^{cl}_{\mathfrak{g}}(O)$ 的描述。若把 $\mathfrak{z}^{cl}_{\mathfrak{g}}(O)$ 与 $\mathfrak{g}^{*}\otimes\omega_{O}$ 上 $G(O)$-不变多项式的空间等同（见 2.4.1），则 $\overline{\mathcal{L}}_{n}$ 的主象征是多项式 $\ell_{n}: \mathfrak{g}^{*}\otimes\omega_{O} \to \mathbb{C}$，由 $\ell_{n}(\eta) = \frac{1}{2}\operatorname{Res}(\eta, \eta)L_{n}$ 定义。这里 $(\eta, \eta) \in \omega_{O}^{\otimes 2}$，$L_{n} \in \omega_{K}^{\otimes(-1)}$，$(\eta, \eta)L_{n} \in \omega_{K}$，故留数有意义。显然映射 $\mathbb{C}[\ell_{-2}, \ell_{-3}, \ldots] \to \mathfrak{z}^{cl}_{\mathfrak{g}}(O)$ 是单射；当 $\mathfrak{g} = \mathfrak{sl}_{2}$ 时它是同构。

对 $\mathfrak{g} = \mathfrak{sl}_{2}$，Feigin--Frenkel 同构是唯一的 $\operatorname{Der} O$-等变同构 $A_{{}^{L}\mathfrak{g}}(O) \xrightarrow{\sim} \mathfrak{z}_{\mathfrak{g}}(O)$。$\operatorname{Spec} O$ 上的 $\mathfrak{sl}_{2}$-oper 可表示成联络 $\frac{d}{dt} + \begin{pmatrix} 0 & u \\ 1 & 0 \end{pmatrix}$，$u = u(t) = u_{0} + u_{1}t + \ldots$；也可表示成 Sturm--Liouville 算子 $(\frac{d}{dt})^{2} - u(t): \omega_{O}^{-1/2} \to \omega_{O}^{3/2}$。有 $A_{\mathfrak{sl}_{2}}(O) = \mathbb{C}[u_{0}, u_{1}, \ldots]$，而 Feigin--Frenkel 同构把 $u_{j}$ 映到 $-2\overline{\mathcal{L}}_{-2-j}$。

对任意半单 $\mathfrak{g}$，我们在 3.5.6 中给出了 $A_{{}^{L}\mathfrak{g}}(O)$ 作为带 $\operatorname{Der} O$ 作用的代数的描述；见（64）--（68）。利用 Feigin--Frenkel 同构的 $\operatorname{Der} O$-等变性可见：若 $\mathfrak{g}$ 单，则 $\overline{\mathcal{L}}_{-2-j} \in \mathfrak{z}_{\mathfrak{g}}(O)$ 对应 $cu_{1j} \in A_{{}^{L}\mathfrak{g}}(O)$，其中 $c = -(\dim\mathfrak{g})/6$ (???)。
% [Tutor Note: 原书 p.97 括注「(???)」为作者自记之疑问，且下标印作「1j」，底本如此，保留原样。]

\bdp{3.6.19} 考虑真空模 $\operatorname{Vac}_{\lambda} := \operatorname{Vac}_{A}/(h-\lambda)\operatorname{Vac}_{A}$，其中 $\operatorname{Vac}_{A}$ 是 $A$ 模去由 $\mathfrak{g}\otimes O$ 生成的闭左理想所得的商。2.9.3 中我们提到，当 $\lambda \ne 0$ 时 $\operatorname{End}_{A}\operatorname{Vac}_{\lambda} = \mathbb{C}$。下述证明

\SourcePage{98}
系由 E. Frenkel 告知我们。正如 2.9.3--2.9.5 所解释的，任何自同态 $f: \operatorname{Vac}_{\lambda} \to \operatorname{Vac}_{\lambda}$ 都来自 $A/(h-\lambda)A$ 的某个中心元素 $z$。事实上，当 $\lambda \ne 0$ 时 $A/(h-\lambda)A$ 的中心等于 $\mathbb{C}$；但我们不去证明这一点，转而注意 $[\widetilde{\mathcal{L}}_{0}, z] = 0$，从而 $L_{0}(z) = 0$（见（87））。于是 $[L_{0}, f] = 0$，其中 $L_{0}$ 看作 $\operatorname{Vac}_{\lambda}$ 中的算子。因此 $f$ 保持空间 $\operatorname{Ker}(L_{0}: \operatorname{Vac}_{\lambda} \to \operatorname{Vac}_{\lambda})$，该空间由真空向量生成。由于 $A$-模 $\operatorname{Vac}_{\lambda}$ 由这个空间生成，$f$ 是纯量算子。

\bdsec{3.7. 中心与 Gelfand--Dikii 括号}

\bdp{3.7.1} 令 $Y := \operatorname{Spec} O$，$Y' := \operatorname{Spec} K$，其中按惯例 $O = \mathbb{C}[[t]]$，$K = \mathbb{C}((t))$。设 $A$ 是（交换的）$\operatorname{Aut} O$-代数。则对任意光滑曲线 $X$ 可得 $\mathcal{D}_{X}$-代数 $A_{X}$（见 2.6.5）。$Y$ 与 $Y'$ 虽不是字面意义上的曲线，2.6.5 的构造对它们仍适用（只需作一点修改，见下）。于是得到 $\mathcal{D}_{Y}$-代数 $A_{Y}$ 与 $\mathcal{D}_{Y'}$-代数 $A_{Y'}$，后者是 $A_{Y}$ 到 $Y'$ 的限制。$A_{Y}$ 在原点 $0 \in Y$ 处的纤维等于 $A$。

下面解释一些细节。2.6.5 中 $A_{X}$ 的定义用到某个概形 $X^{\wedge}$。由于 $Y$ 不是字面意义上的曲线，$Y^{\wedge}$ 的定义应作如下修改。记 $\Delta_{n}$ 为对角线 $\Delta \subset \operatorname{Spec} O\widehat{\otimes}O$ 的第 $n$ 个无穷小邻域。态射 $\operatorname{Spec} O\widehat{\otimes}O \to \operatorname{Spec} O\otimes O = Y\times Y$ 诱导嵌入 $\Delta_{n} \hookrightarrow Y\times Y$（若 $n > 0$，则 $\Delta_{n}$ 比 $Y\times Y$ 中对角线 $\Delta$ 的第 $n$ 个无穷小邻域\emph{更小}）。现在，$Y^{\wedge}$ 的 $R$-点只应考虑满足下述性质的 $R$-态射 $\gamma: \operatorname{Spec} R\widehat{\otimes}O \to Y$：对任意 $n$ 存在 $N$，使得 $\gamma$ 诱导的态射 $\operatorname{Spec} O/t^{n}O \times \operatorname{Spec} O/t^{n}O \times \operatorname{Spec} R \to Y\times Y$ 经过 $\Delta_{N}$ 分解（此时可取 $N = 2n-2$）。

\bdp{3.7.2} 有时我们会用到截面
\begin{equation}\tag{94}
Y \to Y^{\wedge},
\end{equation}

\SourcePage{99}
它对应于由下式定义的态射 $\gamma: \operatorname{Spec} O\widehat{\otimes}O \to Y = \operatorname{Spec} O$：
\begin{equation}\tag{95}
\gamma^{*}(t) = t\otimes 1 + 1\otimes t.
\end{equation}
截面（94）给出同构
\begin{equation}\tag{96}
A_{Y} \xrightarrow{\sim} A\otimes O_{Y}.
\end{equation}
当然，（94）与（96）不是典范的：它们依赖于局部参数 $t \in O$ 的选取。

\bdp{3.7.3} 在 3.7.1 的情形中，考虑从 $\mathbb{C}$-代数到集合的函子 $F$，它使 $F(R)$ 为水平 $Y'$-态射 $\operatorname{Spec} R\widehat{\otimes}K \to \operatorname{Spec} A_{Y'}$ 的集合；等价地，为水平 $K$-态射 $H^{0}(Y', A_{Y'}) \to R\widehat{\otimes}K$ 的集合。$F$ 可由一个归纳仿射归纳概形 $S$ 表示（可称之为 $\operatorname{Spec} A_{Y'}$ 的水平截面归纳概形）。事实上，$F$ 是函子 $R \mapsto \operatorname{Hom}(V, R\widehat{\otimes}K)$ 的闭子函子，其中 $V = H^{0}(Y', A_{Y'})$，$\operatorname{Hom}$ 表示 $K$-线性映射的集合。

记 $A_{K}$ 为 $S$ 上正则函数环。这是完备拓扑代数（$A_{K}$ 中对应于 $S$ 的闭子概形的理想构成零的邻域基）。

$A_{K}$ 带群归纳概形 $\operatorname{Aut} K$ 的作用（$\operatorname{Aut} K$ 的 $R$-点是拓扑 $R$-代数 $R\widehat{\otimes}K$ 的自同构）。

$\operatorname{Spec} A_{Y}$ 的水平截面概形典范同构于 $\operatorname{Spec} A$（水平截面 $s: Y \to \operatorname{Spec} A_{Y}$ 对应到 $s(0) \in \operatorname{Spec} A$）。这是 $S = \operatorname{Spec} A_{K}$ 的闭子概形，于是得到典范满态射
\begin{equation}\tag{97}
A_{K} \to A.
\end{equation}
显然它是 $\operatorname{Aut} O$-等变的。

\emph{例.} 设 $A = \mathbb{C}[u_{0}, u_{1}, u_{2}, \ldots]$，$u_{k} = (L_{-1})^{k}u_{0}/k!$，$L_{0}u_{0} = du_{0}$，$d \in \mathbb{Z}$（按惯例 $L_{n} := -t^{n+1}\frac{d}{dt} \in \operatorname{Der} O$）。则 $\mathcal{D}_{Y}$-概形 $\operatorname{Spec} A_{Y}$ 与 $Y$ 上 $d$-微分的 jets 概形

\SourcePage{100}
之间存在显然的同构 $f$。显然 $\operatorname{Aut} O = \operatorname{Aut} Y$ 按函子性作用在这两个概形上。$f$ 关于由 $L_{0}$ 与 $L_{-1}$ 生成的 $\operatorname{Aut} O$ 群归纳概形等变。借助 $f$，我们把 $\operatorname{Spec} A_{Y'}$ 的水平截面与 $Y'$ 上的 $d$-微分等同，即与 $\omega_{Y'}^{\otimes d}$ 的截面等同。$Y'$ 上的 $d$-微分可写成 $\sum_{i}\tilde{u}_{i}t^{i}(dt)^{\otimes d}$，故 $A_{K} = \mathbb{C}[[\ldots, \tilde{u}_{-1}, \tilde{u}_{0}, \tilde{u}_{1}, \ldots]]$，其中
\begin{equation}\tag{98}
\mathbb{C}[[\ldots, \tilde{u}_{-1}, \tilde{u}_{0}, \tilde{u}_{1}, \ldots]] := \varprojlim_{n}\mathbb{C}[\ldots, \tilde{u}_{-1},, \tilde{u}_{0}, \tilde{u}_{1}, \ldots]/(u_{-n}, u_{-n-1}, \ldots).
\end{equation}
% [Tutor Note: 原书 p.100 式（98）中被取极限的环内印有重复逗号「\tilde{u}_{-1},, \tilde{u}_{0}」，且商理想的生成元印作 $u_{-n}$ 而非 $\tilde{u}_{-n}$；底本如此，保留原样。]
显然 $L_{0}\tilde{u}_{k} = (d+k)\tilde{u}_{k}$，$L_{-1}\tilde{u}_{k} = (k+1)\tilde{u}_{k+1}$，而态射（97）把 $\tilde{u}_{k}$ 映到 $u_{k}$（$k \ge 0$），映到 $0$（$k < 0$）。

\bdp{3.7.4} 记 $\mathfrak{z}_{\mathfrak{g}}(K)$ 为 3.7.3 中代数 $A_{K}$ 在 $A = \mathfrak{z}_{\mathfrak{g}}(O)$ 这一特殊情形的特例（$\mathfrak{z}_{\mathfrak{g}}(O)$ 的定义见 2.5.1 或 2.7.2）。我们将定义从 $\mathfrak{z}_{\mathfrak{g}}(K)$ 到完备化扭泛包络代数 $\overline{U}' = \overline{U}'(\mathfrak{g}\otimes K)$ 的中心 $\mathfrak{Z}$ 的典范态射。为此把（34）改写成 $K$-线性映射 $\mathfrak{z}_{\mathfrak{g}}(O)\otimes_{O}K \to \mathfrak{Z}\widehat{\otimes}K$。利用非典范同构 $\mathfrak{z}_{\mathfrak{g}}(O)_{Y} \xrightarrow{\sim} \mathfrak{z}_{\mathfrak{g}}(O)\otimes O_{Y}$（见（96）），得到映射
\begin{equation}\tag{99}
H^{0}(Y', \mathfrak{z}_{\mathfrak{g}}(O)_{Y'}) \to \mathfrak{Z}\widehat{\otimes}K,
\end{equation}
容易看出它是典范的，即与局部参数 $t \in O$ 的选取无关。事实上，（34）是（99）的非典范版本；（34）依赖于 $t$ 的选取，因为（32）中出现 $\zeta + t$，而它正是由（95）定义的非典范截面 $Y' \to Y'^{\wedge}$。

\bdp{3.7.5} \emph{定理.}
\begin{enumerate}[label=(\roman*),leftmargin=2.2em]
  \item 映射（99）是 $K$-代数的水平态射。因此（99）定义了连续态射
  \begin{equation}\tag{100}
  \mathfrak{z}_{\mathfrak{g}}(K) \to \mathfrak{Z}.
  \end{equation}
  \item 复合 $\mathfrak{z}_{\mathfrak{g}}(K) \to \mathfrak{Z} \to \mathfrak{z}_{\mathfrak{g}}(O)$ 就是取 $A = \mathfrak{z}_{\mathfrak{g}}(O)$ 时的态射（97）。
  \item 态射（100）是 $\operatorname{Aut} K$-等变的。
\end{enumerate}
